\begin{explanationbox}
	\textbf{\textcolor{CExplanation}{一句话直觉}}

	在三角形中，要比较两个角的大小，只需比较它们的正弦值或对边长度。

	\medskip
	\textbf{\textcolor{CExplanation}{核心拆解}}
	\begin{description}[style=nextline, leftmargin=2.8em, labelsep=0.5em, itemsep=0.45em]
		\item[\textbf{要点一（三向等价）：}] 角、边、正弦三者两两等价，任选其一即可。
		\item[\textbf{要点二（单调性陷阱）：}] 正弦函数在$(0,\pi)$上并非单调，但在三角形内角限制下大角对应正弦值更大。
		\item[\textbf{要点三（定理桥梁）：}] 正弦定理将边与正弦联系起来，大边对大角将边与角联系起来，从而形成闭环。
	\end{description}

	\medskip
	\textbf{\textcolor{CExplanation}{几何本质（边角互化）}}

	大边对大角是几何事实，正弦定理则是代数桥梁。两者结合给出一个简捷的边角比较工具。

	\medskip
	\textbf{\textcolor{CExplanation}{代数意义（比较转化）}}

	将几何中角的大小比较转化为正弦值或边长的数值比较，方便运算和推理。

	\medskip
	\textbf{\textcolor{CExplanation}{考点价值（选择填空利器）}}
	\begin{itemize}[leftmargin=*, itemsep=0.5em, topsep=0.4em]
		\item \textbf{考法一（边角互化）：} 已知边的关系，判断角的大小或正弦值大小；反之亦然。
		\item \textbf{考法二（不等式证明）：} 在证明边角不等关系时，可直接套用等价链。
		\item \textbf{考法三（选项排除）：} 利用该结论可快速排除选择题中矛盾的选项。
	\end{itemize}

	\medskip
	\textbf{\textcolor{CExplanation}{顿悟点}}

	在三角形范围内，正弦函数实际上是“局部单调”的——不要被它先增后减的整体图像迷惑。

	\medskip
	\textbf{\textcolor{CExplanation}{使用场景}}
	\begin{itemize}[leftmargin=*, itemsep=0.5em, topsep=0.4em]
		\item \textbf{场景一（直接比较）：} 已知两角或两边大小，判断对应边或角的大小。
		\item \textbf{场景二（解三角形验证）：} 在解三角形时，用该结论检验所得解的合理性（如判定是否增根）。
		\item \textbf{场景三（参数范围）：} 利用边角关系确定某变量的取值范围。
	\end{itemize}
\end{explanationbox}
